\section{Empirical Results}

\label{sec:results}
% ══════════════════════════════════════════════════════════════════════════════

\subsection{Operational Summary}

\begin{table}[H]
\centering
\caption{Zoo DAO operational statistics (18-month period).}
\label{tab:ops}
\begin{tabular}{lc}
\toprule
\textbf{Metric} & \textbf{Value} \\
\midrule
Total proposals submitted & 287 \\
Proposals funded & 143 (49.8\%) \\
Countries represented & 34 \\
Total funding distributed & \$4.7M equivalent \\
Mean project size & \$32{,}900 \\
Unique voters (per round) & 847 $\pm$ 214 \\
Mean review turnaround & 8.4 days \\
Researcher satisfaction & 4.3/5.0 \\
Published outputs & 89 papers, 34 datasets, 21 tools \\
\bottomrule
\end{tabular}
\end{table}

\subsection{Funding Allocation Quality}

\begin{table}[H]
\centering
\caption{Funding allocation quality: QV vs.\ simulated 1T1V.}
\label{tab:allocation}
\begin{tabular}{lcc}
\toprule
\textbf{Metric} & \textbf{QV} & \textbf{1T1V (simulated)} \\
\midrule
Correlation with citation impact & 0.67 & 0.47 \\
Gini coefficient of funding & 0.34 & 0.71 \\
Projects from Global South (\%) & 38\% & 12\% \\
Novel research topics funded (\%) & 44\% & 27\% \\
Post-funding completion rate & 87\% & 82\% \\
\bottomrule
\end{tabular}
\end{table}

QV allocation shows 41\% higher correlation with citation impact, lower
Gini coefficient, triple the Global South funding rate, and more novel
topics funded.

\subsection{Peer Review Performance}

\begin{table}[H]
\centering
\caption{Peer review comparison.}
\label{tab:review_compare}
\begin{tabular}{lcc}
\toprule
\textbf{Metric} & \textbf{Zoo DAO} & \textbf{Traditional} \\
\midrule
Median review time & 8.4 days & 97 days \\
Review completion rate & 94\% & 68\% \\
Author satisfaction & 4.3/5.0 & 3.1/5.0 \\
Review length (words) & 1{,}247 & 892 \\
Actionable suggestions per review & 4.8 & 2.1 \\
\bottomrule
\end{tabular}
\end{table}

Staking produces 10x faster, more complete, and higher-quality reviews.

\subsection{Retroactive Funding Impact}

After three rounds, 47 public goods received awards:
\begin{itemize}
  \item 82\% of funded tools saw increased adoption post-award.
  \item 67\% of funded datasets received new citations within 3 months.
  \item Recipients were 2.4x more likely to contribute subsequent
    public goods.
\end{itemize}


% ══════════════════════════════════════════════════════════════════════════════
